% SHARD-ID: CA-38-QUBIT-LOCAL-ALGEBRA-EQUATION-FRAMEWORK
% SHARD-TITLE: Qubit Local-Algebra Equation Framework
% SHARD-SUMMARY: Fixes the finite Pauli-word calculus used to turn qubit lattice symmetry residuals into explicit coefficient equations.
% SHARD-SUMMARY: Defines the coboundary quotient for infinite-chain bulk densities, so residuals may vanish modulo telescoping boundary terms.
% SHARD-SUMMARY: Separates checked finite algebra from continuum Poincare or Virasoro claims.
% SHARD-KEYWORDS: qubit, Pauli words, local algebra, coboundary, coefficient equations, infinite chain

\section{Qubit Local-Algebra Equation Framework}
\label{sec:qubit-local-algebra-equation-framework}

\subsection{Status, Sources, and Dependencies}

\textbf{Status: checked local framework.}  This shard supplies the algebraic
language used by CA-39--CA-44.  It does not assert that a solution has
continuum Poincare or Virasoro symmetry.

Dependencies: CA-11, CA-12, CA-15, CA-34--CA-35, CONVENTIONS.md (m), (n).

% Source: references/text/CFTFromLatticeFermions.txt:80--100 -- local finite
%   support Hamiltonian terms and observable algebra.
% Source: references/text/CFTFromLatticeFermions.txt:1033--1038 -- Pauli
%   matrices in the qubit setting.
% Source: literature/md/2010.11121/2010.11121.md:204--222 -- local algebras
%   and quasi-local algebra.
% Checked: src/QubitPauliLattice.jl, src/QubitPauliResiduals.jl, and
%   test/runtests.jl, testset "qubit Pauli nearest-neighbour current
%   obstructions".

\subsection{Pauli Word Algebra}

Use zero-based Pauli labels \(a\in\{0,1,2,3\}\), with
\(\sigma^0=I\).  For a finite interval, a Pauli word is
\[
  P_{\alpha_0\ldots\alpha_{\ell-1}}
  =
  \sigma^{\alpha_0}\otimes\cdots\otimes\sigma^{\alpha_{\ell-1}}.
\]
The local product is encoded by constants \(M_{ab}^{c}\):
\[
  \sigma^a\sigma^b=\sum_{c=0}^3 M_{ab}^{c}\sigma^c,
\]
where
\[
  M_{0b}^{c}=\delta_{bc},\qquad
  M_{a0}^{c}=\delta_{ac},
\]
and, for \(a,b\geq1\),
\[
  M_{ab}^{0}=\delta_{ab},\qquad
  M_{ab}^{c}=i\epsilon_{abc}\quad(c\geq1).
\]
Thus for two \(\ell\)-site coefficient tensors \(A,B\),
\[
  (AB)_\gamma
  =
  \sum_{\alpha,\beta}
  A_\alpha B_\beta
  \prod_{r=0}^{\ell-1}M_{\alpha_r\beta_r}^{\gamma_r}.
\]
The commutator coefficient is therefore the finite expression
\[
  (i[A,B])_\gamma
  =
  i\sum_{\alpha,\beta}A_\alpha B_\beta
  \left(
  \prod_rM_{\alpha_r\beta_r}^{\gamma_r}
  -
  \prod_rM_{\beta_r\alpha_r}^{\gamma_r}
  \right).
\]
When two local densities have different supports, first embed both tensors into
their common finite window by padding with Pauli label \(0\), then apply this
formula.

\subsection{Bulk Equality Modulo Divergence}

Let \(\tau\) be the unit shift of the infinite chain.  A finite density \(R\)
represents a zero bulk formal sum if it is a coboundary:
\[
  R=u-\tau(u).
\]
If \(u\) has \(n\) Pauli indices, the coefficient form fixed in CONVENTIONS.md
(n) is
\[
  (D u)_{a_0\ldots a_n}
  =
  u_{a_0\ldots a_{n-1}}\delta_{a_n0}
  -
  \delta_{a_00}u_{a_1\ldots a_n}.
  \tag{38.1}
\]
Consequently a local residual \(R(h)\) passes the bulk equation precisely when
there exists a finite tensor \(u\) such that
\[
  R_{a_0\ldots a_n}(h)-(D u)_{a_0\ldots a_n}=0
  \qquad
  \text{for every Pauli word }a_0\ldots a_n .
  \tag{38.2}
\]
The scalar coefficient of \(u\) lies in the kernel of \(D\), so it is a gauge
variable.  A coefficientwise condition \(R(h)=0\) is allowed as a stronger
filter, but it is not the default infinite-chain equation.

\subsection{What ``Necessary'' Means Here}

Every equation in the next shards is necessary only for the named generator
route: the fixed Pauli density, the fixed first-moment boost ansatz, and the
fixed coboundary quotient.  Solving these equations is not sufficient for
continuum symmetry.  Failing one of them is still useful: it excludes the
Hamiltonian from this route before any low-energy sector, scaling map, or
state-space analysis is attempted.

